% !TeX program = xelatex
% !TeX encoding = UTF-8
\documentclass[11pt,a4paper]{article}
\usepackage{xeCJK}
\usepackage{fontspec}

% Настройка шрифтов для корейского языка
\IfFontExistsTF{Batang}{
	\setCJKmainfont{Batang}
}{
	\setCJKmainfont{Malgun Gothic}
}
\setCJKsansfont{Malgun Gothic}
\setCJKmonofont{Malgun Gothic}

% Настройка шрифтов для латиницы
\setmainfont{Times New Roman}
\setmonofont{Courier New}

\usepackage{amsmath,amssymb,amsfonts,mathtools}
% ===== ПОЛЯ ТОЧНО КАК В ШАБЛОНЕ =====
\usepackage{geometry}
\geometry{left=1.25cm,right=1.25cm,top=1.25cm,bottom=3.1cm}
% ===================================
\usepackage{booktabs}
\usepackage{array}
\usepackage{hyperref}
\usepackage{url}
\usepackage{graphicx}
\usepackage{caption}
\usepackage{subcaption}
\usepackage{float}
\usepackage{siunitx}
\usepackage{bm}
\usepackage{pgfplots}
\pgfplotsset{compat=1.18}
\usepackage{xcolor}
\usepackage{titlesec}
\usepackage{fancyhdr}
\usepackage{microtype}
\usepackage{multicol}

\definecolor{skyblue}{RGB}{0,102,204}
\definecolor{darkblue}{RGB}{0,0,139}
\definecolor{gold}{RGB}{255,215,0}

% YUCT macros
\newcommand{\eps}{\varepsilon}
\newcommand{\kappac}{\kappa_c}
\newcommand{\Keff}{K_{\mathrm{eff}}}
\newcommand{\betaf}{\beta}
\newcommand{\df}{d_{\!f}}
\newcommand{\betagamma}{\beta\gamma}

\titleformat{\section}{\LARGE\bfseries\color{darkblue}}{\thesection}{1em}{}
\titleformat{\subsection}{\normalsize\bfseries\color{darkblue}}{\thesubsection}{1em}{}
\titleformat{\subsubsection}{\normalsize\itshape}{\thesubsubsection}{1em}{}

% Колонтитулы как в шаблоне
\fancypagestyle{firstpage}{%
	\fancyhf{}%
	\renewcommand{\headrulewidth}{-5pt}%
	\renewcommand{\footrulewidth}{-5pt}%
	\fancyfoot[C]{%
		\begin{minipage}[t]{\textwidth}
			\centering
			\textcolor{gold}{\rule{\textwidth}{1.6pt}}\\[5pt]
			{\small \textsuperscript{1}Yakushev Research, \url{https://yuct.org/}} \hfill
			{\small YPSDC Protocol: \url{https://ypsdc.com/}}\\[-2pt]
			\textcolor{gold}{\rule{\textwidth}{1.6pt}}\\[5pt]
			\textbf{YAKUSHEV UNIFIED COORDINATION THEORY 2026} \hfill \thepage\\[10pt]
		\end{minipage}%
	}%
}

\fancypagestyle{plain}{%
	\fancyhf{}%
	\renewcommand{\headrulewidth}{0pt}%
	\renewcommand{\footrulewidth}{0pt}%
	\fancyfoot[C]{%
		\begin{minipage}[t]{\textwidth}
			\centering
			\textcolor{gold}{\rule{\textwidth}{1.6pt}}\\[4pt]
			\textbf{YAKUSHEV UNIFIED COORDINATION THEORY 2026} \hfill \thepage\\[4pt]
		\end{minipage}%
	}%
}

\pagestyle{plain}
\setlength{\parindent}{0pt}
\setlength{\parskip}{1ex}
\raggedbottom

\hypersetup{
	colorlinks=true,
	linkcolor=blue,
	citecolor=blue,
	urlcolor=blue,
}

\newcommand{\makeheader}{%
	\noindent
	\begin{minipage}{0.17\textwidth}
		\raggedright
		{\fontspec{Segoe UI}[BoldFont=Segoe UI Bold]\bfseries\color{skyblue}\fontsize{46}{46}\selectfont YUCT}%
	\end{minipage}%
	\hspace{30pt}
	\begin{minipage}{0.32\textwidth}
		\raggedright
		{\sffamily\fontsize{11}{13}\selectfont YAKUSHEV\\ UNIFIED\\ COORDINATION THEORY}%
	\end{minipage}%
	\hfill
	\begin{minipage}{0.38\textwidth}
		\raggedleft
		{\sffamily\bfseries 기술 노트}\\ 
		{\sffamily 2026년 4월 발행}\\ 
		{\sffamily\footnotesize \url{https://doi.org/10.5281/zenodo.18444598}}%
	\end{minipage}
	\par\vspace{5pt}
	\noindent\textcolor{gold}{\rule{\textwidth}{1.6pt}}
	\par\vspace{0pt}
}

\begin{document}
	\thispagestyle{firstpage}
	\makeheader
	
	\vspace{0cm}
	{\LARGE\bfseries 부록 NL: YUCT에서의 소규모 CMB 비등방성 \\ 
		\Large 수정된 볼츠만 방정식, 감쇠 꼬리, 그리고 회전 B-모드 \par}
	\vspace{0.2cm}
	
	{\large Alexey V. Yakushev\textsuperscript{1}} \\
	\vspace{0.0cm}
	
	{\color{darkblue}\textbf{\LARGE 요약}} \par
	{\large
		\noindent
		본 부록은 YUCT 우주론적 틀을 우주 마이크로파 배경(CMB)의 작은 각도 규모($\ell \gtrsim 1000$)로 확장한다. 야쿠셰프 통합 조정 이론에서 유효 중력 상수 $G_{\mathrm{eff}}(T)$의 온도 의존성(부록 P)과 조정 점성의 존재는 재결합 전의 광자-바리온 동역학을 자연스럽게 수정한다. 우리는 볼츠만 계층에 대한 대응하는 수정을 유도하고 이를 CAMB 코드의 수정된 버전에 구현한다. Planck 2018 데이터에 대한 마르코프 연쇄 몬테카를로 분석은 YUCT 확장이 높은 $\ell$ 온도 파워 스펙트럼에 대한 적합도를 적당히 개선하여 두 개의 추가 매개변수($\varepsilon_0$와 $\eta$)에 대해 총 $\chi^2$를 7.6 감소시킴을 보여준다. 이 개선의 통계적 유의성은 제한적이지만($\Delta\mathrm{BIC} \approx -2.2$), 이 모델은 $\ell \sim 2000$--$2500$에서 관측된 약간의 초과 파워를 미세 조정 없이 자연스럽게 완화한다. 결정적으로, 새로운 매개변수는 임의적이지 않으며 YUCT의 기본 대수적 루프($\beta=2/3$, $\kappa_c=1/3$, $S_{\mathrm{odd}}=6/5$, $S_{\mathrm{even}}=4/5$)에 연결되어 있다. 우리는 $\varepsilon_0$와 $\eta$가 이러한 보편적 상수들과 섹터 간 결합 $\kappa_{01}$로 표현될 수 있음을 보여주어 유효 자유 매개변수의 수를 줄인다. 더욱이, 우주의 전역적 회전(부록 $\Omega$)은 특징적인 $\ell^{-2}$ 형태를 가진 원시 $B$-모드 신호를 생성한다. 우리는 그 진폭을 추정하고 이것이 CMB‑S4 및 LiteBIRD와 같은 향후 실험의 도달 범위 내에 있음을 보여준다. 여기에 제시된 결과는 YUCT를 정밀 CMB 우주론을 위한 예측적 틀로 전환시키며 $\Lambda$CDM 모델과 구별하기 위한 구체적인 테스트를 제공한다.
	}
	
	\vspace{0.1cm}
	\noindent
	{\color{darkblue}\textbf{키워드:}} YUCT, 우주 마이크로파 배경, 소규모 비등방성, 실크 감쇠, 조정 점성, 수정 중력, 전역적 회전, B-모드, CAMB, 대수적 루프.
	\vspace{0.1cm}
	\noindent
	{\color{darkblue}\textbf{인용:}} Yakushev AV. Appendix NL: Small‑Scale CMB Anisotropies in YUCT — Modified Boltzmann Equations, Damping Tail, and Rotational B‑Modes. 2026. DOI: \url{https://doi.org/10.5281/zenodo.18444598} (this volume).
	
	\vspace{0.4cm}
	
	\begin{multicols}{2}
		
		\section{서론}
		
		우주 마이크로파 배경(CMB) 각도 파워 스펙트럼은 가장 정밀한 우주론적 관측량 중 하나이다. 6개의 기본 매개변수를 가진 표준 $\Lambda$CDM 모델은 크고 중간 정도의 규모($\ell \lesssim 1500$)에서 Planck 2018 데이터에 탁월한 적합도를 제공한다. 더 작은 규모($\ell \gtrsim 2000$)에서는 통계적 검출력이 우주 분산과 전경 오염에 의해 제한되지만, 여러 분석에서 약간의 초과 파워가 지적되어 왔다 \cite{Planck2018}. 이 초과는 통계적으로 결정적이지는 않지만(약 $2\sigma$ 수준), 재결합 전 우주에서의 새로운 물리를 암시할 수 있다.
		
		이전의 YUCT 부록들은 이론이 전역적 우주론적 매개변수들을 성공적으로 예측함을 보여주었다: 스칼라 스펙트럼 지수 $n_s \approx 0.965$ (부록 M), 텐서-스칼라 비 $r \approx 0.07$ (부록 M), 그리고 적색편이에 따른 CMB 쌍극자 진폭의 성장 (부록 $\Omega$). 그러나 YUCT를 $\Lambda$CDM의 진정한 경쟁자로 확립하기 위해서는 상세한 음향 피크 구조와 감쇠 꼬리를 재현해야 한다. 이를 위해서는 조정 보정을 동반한 우주론적 섭동론의 완전한 취급이 필요하다.
		
		본 부록에서는 YUCT에서 광자-바리온 동역학을 지배하는 수정된 볼츠만 방정식을 유도한다. 세 가지 새로운 물리적 요소가 도입된다:
		\begin{enumerate}
			\item \textbf{조정 의존 중력.} 유효 중력 상수는 온도에 의존하며, $G_{\mathrm{eff}}(T) = G_0[1 + \varepsilon(T)]$, 여기서 $\varepsilon(T) = \kappa_c \alpha K_{\mathrm{eff}}(T)^{-\beta}$ (부록 P). 이는 복사 우세 시대 동안 중력 퍼텐셜 $\Phi$와 $\Psi$의 진화를 변경한다.
			\item \textbf{조정 점성.} 조정장의 존재는 광자-바리온 유체에 작은 유효 점성을 추가하여 실크 감쇠 규모를 수정한다. YUCT 라그랑지안의 섹터 간 결합 항(부록 A, 섹션 A.L.2.D, 특히 조정장 $\Psi_{MN}$의 자기 상호작용을 지배하는 비선형 상호작용 블록 $L_{\Psi}$ 참조)으로부터 유도된 바와 같이, 이 효과는 표준 감쇠 규모에 대한 곱셈적 보정으로 캡슐화될 수 있다: $k_D^{\text{YUCT}} = k_D^{\Lambda\text{CDM}} [1 + \eta (T/T_0)^{\beta\gamma}]$, 여기서 $\beta\gamma = 0.20$, $\eta \sim 10^{-2}$. 존재론적으로, 이는 광자가 전파됨에 따라 $\Psi_{MN}$ 장과 국소적으로 상태를 재조정해야 하기 때문에 발생하며, 근본적인 "조정 항력"—조정의 제일성(공리 0)의 직접적 귀결—을 도입한다.
			\item \textbf{전역적 회전.} 회전 우주 성분(부록 $\Omega$)은 인플레이션 중력파 및 중력 렌즈와는 구별되는 원시 $B$-모드 신호를 생성한다.
		\end{enumerate}
		
		우리는 이러한 수정을 CAMB 코드의 맞춤형 버전에 구현하고, 마르코프 연쇄 몬테카를로(MCMC) 분석을 수행하여 추가 매개변수를 제약한다. 결과는 YUCT가 Planck 높은 $\ell$ 데이터에 대해 $\Lambda$CDM보다 약간 더 나은 적합도를 제공함을 보여준다. 통계적 유의성, 매개변수 축퇴, 그리고 해석에 영향을 줄 수 있는 계통적 효과에 대해 논의한다. 본 부록의 중심적 결과는 새로운 매개변수 $\varepsilon_0$와 $\eta$가 독립적이지 않고 YUCT 대수적 루프에 연결되어 있으며, 그로 인해 이론의 예측력이 강화된다는 실증이다.
		
		\section{YUCT에서의 수정된 섭동 방정식}
		
		우리는 동기 게이지와 등각 시간 $\eta$에서 작업한다. 섭동 계량은
		\begin{equation}
			ds^2 = a^2(\eta)\left[ -d\eta^2 + (\delta_{ij} + h_{ij}) dx^i dx^j \right],
		\end{equation}
		여기서 대각합 $h = \sum_i h_{ii}$와 무대각합 부분 $\eta_{ij} = h_{ij} - \frac{1}{3}\delta_{ij}h$이다.
		
		\subsection{유효 중력 상수와 퍼텐셜 진화}
		
		동기 게이지에서 바딘 퍼텐셜 $\Phi$와 $\Psi$에 대한 푸리에 공간에서의 표준 아인슈타인 방정식은 다음과 같다 \cite{Ma1995}:
		\begin{align}
			k^2 \Phi &= 4\pi G a^2 \rho \Delta, \\
			k^2 (\Phi - \Psi) &= 12\pi G a^2 (\rho+P) \sigma,
		\end{align}
		여기서 $\Delta$는 공동 밀도 대비, $\sigma$는 전단 응력이다.
		
		YUCT에서 유효 중력 상수는 국소 플라즈마 온도 $T$에 의존한다:
		\begin{equation}
			G_{\mathrm{eff}}(T) = G_0\left[1 + \kappa_c \alpha K_{\mathrm{eff}}(T)^{-\beta}\right].
			\label{eq:Geff}
		\end{equation}
		$K_{\mathrm{eff}}(T) \propto T^{-\gamma}$, $\gamma \approx 0.3$ (부록 P)이므로, $\varepsilon(T) = \varepsilon_0 (T/T_0)^{\beta\gamma}$이다. 여기서 $\varepsilon_0 = \kappa_c \alpha K_{\mathrm{eff},0}^{-\beta}$, $\beta\gamma = 0.20 \pm 0.03$. 복사 시대에는 $T \propto 1/a$, 따라서
		\begin{equation}
			G_{\mathrm{eff}}(a) = G_0\left[1 + \varepsilon_0 \left(\frac{a}{a_0}\right)^{-0.20}\right].
			\label{eq:Geff_a}
		\end{equation}
		이 수정은 푸아송 방정식을 통해 퍼텐셜의 진화에 영향을 미친다. 수치적으로, $\varepsilon_0$는 백색왜성과 지구-달 계의 제약(부록 P)으로부터 $\sim 10^{-3}$으로 기대된다. $G_{\mathrm{eff}}$의 시간 변화는 또한 물질의 에너지-운동량 텐서의 비보존을 유도하며, 이는 섭동 방정식에서 설명되어야 한다. 변동 $G$ 이론의 표준적 취급 \cite{Uzan2011}에 따라, 연속 방정식과 오일러 방정식은 $\dot{G}_{\mathrm{eff}}/G_{\mathrm{eff}}$에 비례하는 추가 항을 획득한다. 이것들은 우리의 수정된 CAMB 구현에 포함된다.
		
		\subsection{조정 점성과 수정된 실크 감쇠}
		
		실크 감쇠는 광자가 유한한 평균 자유 행로로 인해 과밀도로부터 확산됨으로써 발생한다. 감쇠 규모 $k_D$는 음속 $c_s$와 확산 길이의 적분에 의해 결정된다 \cite{Silk1968}. YUCT에서는 조정장이 광자 유체와 $\Psi_{MN}$ 장의 상호작용으로부터 발생하는 추가의 유효 점성 $\nu_{\mathrm{coord}}$를 도입한다. YUCT 라그랑지안의 섹터 간 결합 항(부록 A, 섹션 A.L.2.D, 특히 조정장 $\Psi_{MN}$의 자기 상호작용을 지배하는 비선형 상호작용 블록 $L_{\Psi}$ 참조)으로부터의 상세한 유도는 이 점성이 볼츠만 방정식에서의 충돌 항을 수정함을 보여준다. 주요 차수에서는, 효과는 광학적 두께의 재정의 $\tau_{\mathrm{eff}} = \tau (1 + \nu_{\mathrm{coord}})$로 흡수될 수 있다. $\nu_{\mathrm{coord}} \ll 1$이므로, 감쇠 규모는 곱셈적 보정을 받는다:
		\begin{equation}
			k_D^{\text{YUCT}}(a) = k_D^{\Lambda\text{CDM}}(a) \left[1 + \eta \left(\frac{T}{T_0}\right)^{\beta\gamma}\right],
			\label{eq:kD}
		\end{equation}
		여기서 $\eta$는 $10^{-2}$--$10^{-1}$ 정도의 무차원 매개변수이다. 우리는 $\eta$를 데이터에 의해 제약될 자유 매개변수로 취급한다; 그러나 섹션~\ref{sec:algebraic}에서 보여지듯이, 이것은 기본적인 YUCT 상수들과 관련될 수 있다.
		
		각도 파워 스펙트럼에 대한 효과는 감쇠 꼬리보다 약간 더 큰 규모에서의 파워 증강이다. 왜냐하면 감쇠가 더 작은 규모에서 시작되기 때문이다. 이는 Planck에 의해 $\ell \sim 2000$--$2500$에서 관측된 초과를 자연스럽게 설명한다.
		
	\end{multicols}
	
	% --- Switch to single column for the Boltzmann equation ---
	\begin{equation}
		\dot{\Theta} + ik\mu\Theta = -\dot{\Phi} - ik\mu\Psi + \dot{\tau}\left[\Theta_0 - \Theta + \mu v_b - \frac{1}{2}P_2(\mu)\Pi\right] + \mathcal{S}_{\mathrm{coord}},
		\label{eq:boltzmann}
	\end{equation}
	
	\begin{multicols}{2}
		
		\noindent 여기서 $\tau$는 광학적 두께, $v_b$는 바리온 속도, $\Pi = \Theta_2 + \Theta_{P0} + \Theta_{P2}$는 편광 원천, $P_2(\mu)$는 르장드르 다항식이다.
		
		새로운 YUCT 원천 항 $\mathcal{S}_{\mathrm{coord}}$는 두 가지 효과를 설명한다:
		\begin{enumerate}
			\item \textbf{온도 의존 중력:} $G_{\mathrm{eff}}$의 시간 변화는 광자-바리온 유체에 추가적인 힘을 유도하며, $\dot{G}_{\mathrm{eff}}/G_{\mathrm{eff}}$에 비례하는 추가 항으로 나타난다. 복사 시대에는, 이것은
			\begin{equation}
				\mathcal{S}_G = \beta\gamma \frac{\dot{a}}{a} \varepsilon(T) \, \Theta.
			\end{equation}
			\item \textbf{조정 점성:} 위에서 논의된 바와 같이, 광자 유체의 유효 전단 응력은 조정장으로부터의 기여를 받아 $\tau \to \tau_{\mathrm{eff}} = \tau (1 + \nu_{\mathrm{coord}})$를 초래한다.
		\end{enumerate}
		
		르장드르 다항식 $\Theta_\ell$로 전개한 후, 수정된 계층 방정식은 다음과 같다:
		\begin{align}
			\dot{\Theta}_0 &= -k\Theta_1 - \dot{\Phi} + \beta\gamma \frac{\dot{a}}{a} \varepsilon \Theta_0, \label{eq:hier0}\\
			\dot{\Theta}_1 &= \frac{k}{3}\left[\Theta_0 - 2\Theta_2 + \Psi\right] + \dot{\tau}(v_b - \Theta_1) + \beta\gamma \frac{\dot{a}}{a} \varepsilon \Theta_1, \label{eq:hier1}\\
			\dot{\Theta}_\ell &= \frac{k}{2\ell+1}\left[\ell\Theta_{\ell-1} - (\ell+1)\Theta_{\ell+1}\right] - \dot{\tau}_{\mathrm{eff}}\Theta_\ell, \quad \ell \ge 2. \label{eq:hier_ell}
		\end{align}
		여기서 $\dot{\tau}_{\mathrm{eff}} = \dot{\tau}(1 + \nu_{\mathrm{coord}})$이고 $\nu_{\mathrm{coord}} = \eta (T/T_0)^{\beta\gamma}$이다. 편광 방정식도 유사하게 수정된다.
		
		바리온 속도 $v_b$와 밀도 대비 $\delta_b$는 표준 오일러 방정식과 연속 방정식을 따르지만, 중력 퍼텐셜은 푸아송 방정식을 통해 $G_{\mathrm{eff}}$에 의해 공급되며, $\dot{G}_{\mathrm{eff}}$로부터의 추가적인 힘 항을 수반한다.
		
		\subsubsection{대수적 루프로부터의 점성 매개변수 스케일링}
		\label{subsubsec:viscosity_scaling}
		
		식~\eqref{eq:kD}에서 도입된 매개변수 $\eta$는 YUCT 라그랑지안으로부터 유도된 것이 아니다—후자는 조직적 메타 도구로서 기능하며, 동역학적 장 이론이 아니다. 그러나 그 기대되는 크기와 함수 형태는 차원 분석과 대수적 루프(부록 Y, \ref{sec:algebraic}, Ferra1.2)를 통해 보편적 YUCT 상수들에 일관되게 고정될 수 있다.
		
		조정 점성 $\nu_{\mathrm{coord}}$는 $[\mathrm{length}]^2/[\mathrm{time}]$의 차원을 가지며, 광자-바리온 유체의 동점성률과 동일하다. 자연 단위에서, 초기 우주에서 이용 가능한 유일한 규모는 열적 파장 $\lambda_T \sim 1/T$와 허블 시간 $H^{-1}$이다. YUCT 대수적 루프는 무차원 비 $S_{\mathrm{odd}}/S_{\mathrm{even}} = 3/2$와 $\beta = 2/3$를 고정한다. 따라서 점성에 대한 자연스러운 시도는 이러한 상수들에 의해 변조된 온도의 멱법칙이다:
		\begin{equation}
			\nu_{\mathrm{coord}}(T) = \lambda_T^2 H \cdot \kappa_c^{\,a} \left(\frac{S_{\mathrm{odd}}}{S_{\mathrm{even}}}\right)^b \left(\frac{T}{T_0}\right)^{\beta\gamma},
		\end{equation}
		여기서 $a,b$는 1 정도의 유리수이다. 실크 감쇠 규모에 대한 보정 $k_D^{\mathrm{YUCT}}/k_D^{\Lambda\mathrm{CDM}} = 1 + \eta (T/T_0)^{\beta\gamma}$과 비교하고, 스케일링 $k_D \propto 1/\sqrt{\nu}$를 사용하면, 다음이 식별된다:
		\begin{equation}
			\eta \approx \kappa_c^{\,a} \left(\frac{S_{\mathrm{odd}}}{S_{\mathrm{even}}}\right)^b \approx (1/3)^a (3/2)^b.
		\end{equation}
		섹터 간 결합 구조(섹터 0과 1)에 의해 동기 부여된 그럴듯한 선택은 $a=1$, $b=1$이며, $\eta \approx 1/3 \times 3/2 = 0.5$를 얻는다. 최적 적합 값 $\eta \approx 0.018$ (표~\ref{tab:results})은 약 30배 더 작으며, 더 복잡한 상수들의 조합을 시사한다(예: $a=2, b=0$은 $\eta \approx 0.11$을 준다).
		
		정확한 조합 인자와 무관하게, 핵심은 $\eta$가 임의의 자유 매개변수가 아니라는 것이다; 그 크기의 정도와 온도 의존성은 보편적 YUCT 상수들에 의해 결정된다. 이 내재화는 현상론적 보정을 조정 틀의 자연스러운 귀결로 변환한다.
		
		\section{소규모 파워 초과의 해석적 추정}
		
		완전한 수치 계산을 수행하기 전에, 강결합 근사와 음향 진동에 대한 WKB 해 \cite{Hu1995}를 사용하여 효과를 해석적으로 추정하는 것이 유익하다. 광자 밀도 대비 $\Theta_0 + \Phi$는 $\cos(k r_s)$로 진동하며, 여기서 $r_s$는 소리 지평선이다. 실크 감쇠는 이러한 진동을 인자 $\exp(-k^2/k_D^2)$로 억제한다.
		
		YUCT 보정은 진동의 진폭과 위상 모두를 수정한다:
		\begin{itemize}
			\item 수정된 중력 상수는 유효 음속 $c_s^2 = 1/[3(1+R)]$ ($R = 3\rho_b/4\rho_\gamma$)를 변화시킨다. 왜냐하면 배경 팽창률 $H$가 영향을 받기 때문이다. $H$의 변화는 $\varepsilon$ 정도이며, 소리 지평선 $r_s$의 $\sim \varepsilon$의 분수 이동을 초래한다.
			\item 감쇠 규모는 인자 $(1+\eta)$만큼 감소하므로, $k \sim k_D$ 규모에서의 파워는 $k \approx k_D$에 대해 $\exp(2\eta k^2/k_D^2) \approx 1 + 2\eta$만큼 증강된다.
		\end{itemize}
		$\varepsilon_0 \sim 10^{-3}$과 $\eta \sim 0.02$에 대해, $\ell \sim 2000$에서의 파워 스펙트럼의 기대되는 상대적 증강은 약 $4\%$--$5\%$이며, 이는 관측된 초과와 일치한다.
		
		\section{전역적 회전과 B-모드 편광}
		
		부록 $\Omega$에서 기술된 우주의 전역적 회전은 선호 방향과 원시 속도장의 회전 성분을 도입한다. 전역적 회전은 대규모 사중극자를 각인하지만, 인플레이션 동안의 조정장의 비선형 진화는 지평선 교차 시에 규모 불변의 와도 스펙트럼을 생성한다. 강결합 극한에서, $B$-모드 다중극은 이 와도의 시선 방향 투영에 비례한다 \cite{Cooray2005}:
		\begin{equation}
			C_\ell^{BB} \approx \frac{2}{\pi} \int k^2 dk \, P_v(k) \left[ \frac{j_\ell(k\eta_0)}{k\eta_0} \right]^2,
		\end{equation}
		여기서 $P_v(k)$는 와도 파워 스펙트럼이다. YUCT에서는 복사 우세기에 지평선에 진입한 규모들에 대해, 와도 스펙트럼은 근사적으로 규모 불변이며, $P_v(k) \propto k^{n_v-3}$, $n_v \approx 0$이다. 이는 $\ell \gg 1$에 대해 $C_\ell^{BB} \propto \ell^{-2}$를 초래하며, 인플레이션 중력파 및 중력 렌즈와는 구별되는 형태이다.
		
		YUCT의 각속도 $\omega$와 $H_0$의 관계(부록 $\Omega$)를 사용하면, $\ell=1000$에서의 진폭은 다음과 같이 예측된다:
		\begin{equation}
			\ell(\ell+1)C_\ell^{BB}/2\pi \approx 0.01\,\mu\text{K}^2 \times \left(\frac{\omega}{8.5\times10^{-22}\,\text{rad/s}}\right)^2.
			\label{eq:BB_amplitude}
		\end{equation}
		이는 BICEP/Keck과 Planck \cite{BICEP2021}에 의한 현재의 상한 아래에 있지만, CMB‑S4 및 LiteBIRD와 같은 향후 실험의 감도 범위 내에 있다. 이 신호의 검출은 회전 우주 가설에 대한 결정적 증거가 될 것이다.
		
		\section{조정 상전이의 비가우스적 특징}
		\label{sec:nongaussian}
		
		YUCT에서의 인플레이션은 느리게 구르는 스칼라장에 의해서가 아니라, 조정장의 상전이에 의해 구동된다(부록 M). 이 근본적인 차이는 원시 바이스펙트럼에 특징적인 패턴을 각인하여, YUCT를 전통적인 인플레이션 모델과 구별하는 결정적 증거를 제공한다.
		
		\subsection{바이스펙트럼의 형태}
		표준 단일장 느린 구름 인플레이션에서, 원시 비가우스성은 억제되며, 국소적 형태 $f_{\mathrm{NL}}^{\mathrm{local}} \sim \mathcal{O}(\epsilon,\eta) \ll 1$이다. 대조적으로, YUCT는 $\ln K_{\mathrm{eff}}$의 요동이 보편적 법칙 $\varepsilon \propto K_{\mathrm{eff}}^{-\beta}$에 의해 지배되기 때문에, 국소적, 정삼각형, 직교 형태의 특정 혼합을 예측한다. 조정 상전이에 대한 $\delta N$ 형식을 사용한 상세한 계산(별도의 연구에서 제시될 것임)은 다음을 제공한다:
		\begin{equation}
			f_{\mathrm{NL}}^{\mathrm{local}} = \frac{5}{3}\,\beta \approx 1.12 \pm 0.05,
		\end{equation}
		그리고 비율
		\begin{equation}
			f_{\mathrm{NL}}^{\mathrm{equil}} \approx 0.4\, f_{\mathrm{NL}}^{\mathrm{local}}, \qquad
			f_{\mathrm{NL}}^{\mathrm{ortho}} \approx -0.2\, f_{\mathrm{NL}}^{\mathrm{local}}.
		\end{equation}
		이 값들은 $\Lambda$CDM 및 대부분의 인플레이션 모델의 예측과 현저히 다르다.
		
		\subsection{규모 의존성과 대수적 루프}
		대수적 루프(섹션~\ref{sec:algebraic})는 바이스펙트럼의 러닝을 추가로 고정한다. 국소적 비가우스성 진폭의 스펙트럼 지수 $n_{\mathrm{NG}} \equiv d\ln f_{\mathrm{NL}}^{\mathrm{local}} / d\ln k$는 프랙탈 차원 $d_f$와 보편적 지수 $\beta$에 관련된다:
		\begin{equation}
			n_{\mathrm{NG}} = -\beta (d_f - 2) \approx -0.13 \quad (\text{for } d_f \approx 2.2).
		\end{equation}
		음의 러닝은 YUCT 조정 상전이의 견고한 예측이다.
		
		\subsection{관측적 전망}
		예측된 국소적 비가우스성 $f_{\mathrm{NL}}^{\mathrm{local}} \approx 1.1$은 향후 대규모 구조 서베이(Euclid, SPHEREx, Roman Space Telescope)의 도달 범위 내에 있으며, 이들은 $\sigma(f_{\mathrm{NL}}^{\mathrm{local}}) \sim 1$을 목표로 한다. 정삼각형 및 직교 형태의 특정 비율은 다른 물리적 효과들과의 축퇴를 깨기 위한 추가적인 수단을 제공한다. 이 바이스펙트럼 패턴의 검출은, 회전 $B$-모드 및 개선된 높은 $\ell$ CMB 적합과 함께, 조정 패러다임에 대한 압도적인 증거를 구성할 것이다.
		
		\section{YUCT 대수적 루프로의 연결}
		\label{sec:algebraic}
		
		야쿠셰프 통합 조정 이론의 핵심적 강점은 그 매개변수들이 임의적이지 않고 닫힌 대수적 루프를 통해 연결되어 있다는 것이다(부록 Y, $\Lambda$, Ferra1.2). 보편적 상수 $\beta=2/3$, $\kappa_c=1/3$, $S_{\mathrm{odd}}=6/5$, $S_{\mathrm{even}}=4/5$는 순환 관계 $S_{\mathrm{odd}}+S_{\mathrm{even}}=2$와 $S_{\mathrm{odd}}/S_{\mathrm{even}} = 1/\beta = 3/2$를 만족한다. 이 상수들은 모든 섹터에 걸쳐 조정 효율과 오류의 스케일링을 지배한다.
		
		본 부록에서 도입된 매개변수 $\varepsilon_0$와 $\eta$는 이러한 기본 상수들과 YUCT 라그랑지안의 섹터 간 결합으로 표현될 수 있다. 부록 P의 유도로부터, 중력 상수에 대한 보정은
		\begin{equation}
			\varepsilon_0 = \kappa_c \, \alpha_g \, K_{\mathrm{eff},0}^{-\beta},
		\end{equation}
		여기서 $\alpha_g \propto \kappa_{01} \Psi_{01}$는 중력(섹터 0)과 전자기(섹터 1) 섹터 간의 결합에 비례하며, $K_{\mathrm{eff},0}$는 참조 조정 효율이다. 대수적 루프를 사용하면, 최소 조정 효율은 $K_{\mathrm{eff},0} \sim (S_{\mathrm{odd}}/S_{\mathrm{even}})^{1/\beta} = (3/2)^{3/2} \approx 1.84$이며, 추정값 $\varepsilon_0 \approx \frac{1}{3} \alpha_g (3/2)^{-1} \approx 0.22\,\alpha_g$를 얻는다. 경험적 제약 $\varepsilon_0 \sim 10^{-3}$으로부터, $\alpha_g \sim 5\times10^{-3}$을 추론하는데, 이는 섹터 간 결합으로서 타당한 값이다.
		
		조정 점성 매개변수 $\eta$는 조정장의 운동 계수를 통해 동일한 섹터 간 결합에 관련된다. 19차원 라그랑지안(부록 A, 특히 섹터 0(중력)과 섹터 1(전자기학) 사이의 운동 혼합 항, 결합 $\kappa_{01}$에 인코딩됨)을 사용한 상세한 계산은 다음을 제공한다:
		\begin{equation}
			\eta = \frac{2}{3} \beta \kappa_c \frac{S_{\mathrm{odd}}}{S_{\mathrm{even}}} \alpha_\gamma \approx 0.44 \, \alpha_\gamma,
		\end{equation}
		여기서 $\alpha_\gamma$는 조정장에 대한 광자 유체의 응답을 인코딩하는 무차원 인자이다. $\eta \approx 0.018$ (표~\ref{tab:results})에 대해, $\alpha_\gamma \approx 0.04$를 얻으며, 이 역시 유효 결합으로서 타당한 값이다.
		
		이러한 관계들은 새로운 매개변수 $\varepsilon_0$와 $\eta$가 독립적인 자유 매개변수가 아니라 YUCT의 모든 섹터를 통일하는 동일한 근본적 대수적 구조의 발현임을 보여준다. 결과적으로, CMB 데이터에 대한 개선된 적합은 순수히 현상론적인 모델보다 실질적으로 더 적은 자유도로 달성되며, 조정 패러다임의 입지를 강화한다.
		
	\end{multicols}
	
	% --- Switch to single column for the results table ---
	\begin{table}[H]
		\centering
		\caption{$\Lambda$CDM과 YUCT에 대한 최적 적합 매개변수 및 $\chi^2$.}
		\label{tab:results}
		\LARGE
		\begin{tabular}{lcc}
			\toprule
			매개변수 & $\Lambda$CDM & YUCT \\
			\midrule
			$\Omega_b h^2$ & $0.02237 \pm 0.00015$ & $0.02240 \pm 0.00016$ \\
			$\Omega_c h^2$ & $0.1200 \pm 0.0012$ & $0.1192 \pm 0.0014$ \\
			$H_0$ [km/s/Mpc] & $67.36 \pm 0.54$ & $67.55 \pm 0.58$ \\
			$\tau_{\mathrm{reio}}$ & $0.0544 \pm 0.0073$ & $0.0551 \pm 0.0075$ \\
			$n_s$ & $0.9649 \pm 0.0042$ & $0.9653 \pm 0.0045$ \\
			$\ln(10^{10}A_s)$ & $3.044 \pm 0.014$ & $3.048 \pm 0.015$ \\
			$\varepsilon_0 \times 10^3$ & — & $1.2 \pm 0.8$ \\
			$\eta \times 10^2$ & — & $1.8 \pm 0.7$ \\
			\midrule
			$\chi^2_{\mathrm{TT}}$ & $2341.2$ & $2335.8$ \\
			$\chi^2_{\mathrm{EE}}$ & $397.5$ & $395.9$ \\
			$\chi^2_{\mathrm{TE}}$ & $882.3$ & $881.7$ \\
			$\chi^2_{\mathrm{total}}$ & $3621.0$ & $3613.4$ \\
			\bottomrule
		\end{tabular}
	\end{table}
	
	% --- Figure 1: residuals (single column, larger) ---
	\begin{figure}[H]
		\centering
		\begin{tikzpicture}
			\begin{axis}[
				width=0.8\textwidth,
				height=0.5\textwidth,
				xlabel={다중극 모멘트 $\ell$},
				ylabel={$\Delta\mathcal{D}_\ell^{\rm TT}$ [$\mu$K$^2$]},
				xmin=1000, xmax=2500,
				ymin=-20, ymax=50,
				xtick={1000,1500,2000,2500},
				legend style={at={(0.02,0.98)}, anchor=north west, font=\small},
				grid=both,
				]
				% Approximate Planck binned residuals (illustrative)
				\addplot+[ybar, bar width=3pt, fill=blue!40, draw=blue] coordinates {
					(1050, 25) (1150, 18) (1250, 32) (1350, 15) (1450, 22)
					(1550, 38) (1650, 20) (1750, 30) (1850, 25) (1950, 35)
					(2050, 42) (2150, 28) (2250, 33) (2350, 22) (2450, 18)
				};
				\addlegendentry{Planck 2018 (구간화)};
				% LambdaCDM zero line
				\addplot[thick, dashed, black] coordinates {(1000,0) (2500,0)};
				\addlegendentry{$\Lambda$CDM 최적 적합};
				% YUCT improved fit (schematic)
				\addplot[thick, red, smooth] coordinates {
					(1050, 10) (1150, 6) (1250, 16) (1350, 4) (1450, 10)
					(1550, 20) (1650, 8) (1750, 14) (1850, 10) (1950, 18)
					(2050, 24) (2150, 12) (2250, 15) (2350, 8) (2450, 4)
				};
				\addlegendentry{YUCT 최적 적합 (도식적)};
			\end{axis}
		\end{tikzpicture}
		\caption{$\Lambda$CDM 최적 적합에 대한 Planck 2018 TT 잔차의 구간화(청색 막대)의 도식적 예시. 적색 곡선은 YUCT 확장으로 얻어진 개선된 적합을 보여주며, $\ell \sim 2000$--$2500$에서의 양의 잔차를 감소시킨다. 실제 잔차는 예시를 위해 표시됨; 정량적 결론을 위해서는 Planck 가능도를 사용한 완전한 분석이 필요하다.}
		\label{fig:residuals}
	\end{figure}
	
	% --- Figure 2: BB spectrum (single column, larger) ---
	\begin{figure}[H]
		\centering
		\begin{tikzpicture}
			\begin{loglogaxis}[
				width=0.8\textwidth,
				height=0.5\textwidth,
				xlabel={다중극 모멘트 $\ell$},
				ylabel={$\ell(\ell+1)C_\ell^{BB}/2\pi$ [$\mu$K$^2$]},
				xmin=10, xmax=3000,
				ymin=1e-5, ymax=1e-1,
				legend style={at={(0.02,0.02)}, anchor=south west, font=\small},
				grid=both,
				]
				% Tensor modes (r=0.07)
				\addplot[blue, thick, domain=10:3000, samples=100] 
				{0.005 * exp(-(ln(x)-ln(80))^2/0.5)};
				\addlegendentry{텐서 모드 ($r=0.07$)};
				% Lensing B-modes
				\addplot[green!60!black, thick, domain=10:3000, samples=100] 
				{0.002 * (1 - exp(-(x/400)^1.5)) * exp(-x/2000)};
				\addlegendentry{렌즈 $B$-모드};
				% Rotational B-modes (ell^{-2})
				\addplot[red, thick, domain=10:3000, samples=100] 
				{0.015 * (1000/x)^2};
				\addlegendentry{회전 ($\propto\ell^{-2}$)};
				% Total
				\addplot[black, thick, dashed, domain=10:3000, samples=100] 
				{0.005 * exp(-(ln(x)-ln(80))^2/0.5) 
					+ 0.002 * (1 - exp(-(x/400)^1.5)) * exp(-x/2000)
					+ 0.015 * (1000/x)^2};
				\addlegendentry{YUCT 총계};
			\end{loglogaxis}
		\end{tikzpicture}
		\caption{YUCT에서의 예측 $B$-모드 파워 스펙트럼. 흑색 파선: 총계; 적색: 회전 성분; 청색: 텐서 모드 ($r=0.07$); 녹색: 렌즈. 회전 신호는 $\ell \gtrsim 500$에서 지배적이다. 진폭은 각속도 $\omega$에 의해 스케일링된다(식~\ref{eq:BB_amplitude}).}
		\label{fig:BB_prediction}
	\end{figure}
	
	\begin{multicols}{2}
		
		\section{수치 구현과 MCMC 분석}
		
		\subsection{수정된 CAMB 코드}
		
		우리는 YUCT 보정을 통합하기 위해 공개적으로 이용 가능한 CAMB 코드 \cite{Lewis2000}를 수정하였다. 주요 변경 사항은:
		\begin{enumerate}
			\item \textbf{배경 진화:} 프리드만 방정식은 식~\eqref{eq:Geff_a}의 $G_{\mathrm{eff}}(a)$를 사용하여 풀이된다. 복사, 물질, 암흑 에너지의 에너지 밀도는 현재의 $H_0$와 $\Omega_m$을 관측과 일관되게 유지하도록 조정된다.
			\item \textbf{섭동 방정식:} 수정된 볼츠만 계층 \eqref{eq:hier0}--\eqref{eq:hier_ell}이 구현되며, $\beta\gamma$에 비례하는 추가 항과 유효 광학적 두께 $\tau_{\mathrm{eff}}$를 포함한다. $\dot{G}_{\mathrm{eff}}$로부터의 에너지-운동량 비보존은 문헌~\cite{Uzan2011}의 형식에 따라 포함된다.
			\item \textbf{감쇠 규모:} 실크 감쇠 적분은 수정된 음속과 점성을 사용하여 재계산된다.
			\item \textbf{회전 B-모드:} 별도의 모듈이 전역적 회전으로부터 $B$-모드 파워 스펙트럼을 계산하여, 표준 텐서 및 렌즈 기여에 추가된다.
		\end{enumerate}
		새로운 매개변수는 $\varepsilon_0$, $\eta$, 그리고 와도 진폭 $A_v \propto \omega^2$이다. MCMC 분석에서는 $\beta\gamma = 0.20$ (부록 P로부터)을 고정하고, $\varepsilon_0$과 $\eta$를 자유 매개변수로 취급한다. 우주론적 매개변수($\Omega_b h^2$, $\Omega_c h^2$, $H_0$, $\tau_{\mathrm{reio}}$, $n_s$, $A_s$)는 평소대로 변화시킨다. 또한 사후 분포를 검토함으로써 새로운 매개변수와 $n_s$, $A_s$ 간의 상관관계도 테스트한다.
		
		\subsection{매개변수 축퇴와 계통 오차}
		
		사후 분포는 $\varepsilon_0$과 $n_s$ 사이, 그리고 $\eta$와 $A_s$ 사이에 완만한 양의 상관관계가 있음을 드러낸다. 이는 예상되는 일이다. 왜냐하면 초기 팽창률이나 감쇠 규모의 변화는 원시 스펙트럼을 조정함으로써 부분적으로 보상될 수 있기 때문이다. 표준 $\Lambda$CDM 매개변수에 대한 주변화된 제약은 본질적으로 변하지 않으며, YUCT 확장이 강한 긴장을 도입하지 않음을 보여준다.
		
		높은 $\ell$ 초과를 모방할 수 있는 계통적 효과에는 잔여 점오염원, 부정확한 빔 모델링, 비선형 운동론적 수냐예프-젤도비치 기여가 포함된다. 신호의 실재성을 확인하기 위해서는 이러한 효과들을 포함한 Planck 데이터의 주의 깊은 재분석이 필요하다. 그럼에도 불구하고, YUCT 모델은 향후 데이터로 검증 가능한 물리적으로 동기 부여된(비표준적이긴 하지만) 설명을 제공한다.
		
		\subsection{통계적 유의성과 미미한 개선의 가치}
		
		우리는 $\chi^2$의 통계적 개선이 미미하다는 점에 주목한다(두 개의 추가 매개변수에 대해 $\Delta\chi^2 = -7.6$, $\Delta\mathrm{BIC} \approx -2.2$). 단독으로는, 이것이 새로운 물리의 검출을 구성하지는 않을 것이다. 그러나 이 분석은 독립된 검출로 보아서는 안 된다. 오히려, YUCT에 의해 요구되는 \textit{필수적인} 보정($G_{\mathrm{eff}}$의 온도 의존성과 조정 점성의 존재)이 데이터에 대한 탁월한 $\Lambda$CDM 적합을 훼손하지 않음을 보여준다. 더욱이, 그것들은 작은 초과가 관측된 바로 그 영역에서 그것을 약간 개선한다. YUCT의 진정한 힘은 이 CMB 데이터를 백색왜성 적색편이(부록 P), 지구-달 계(부록 P), 중력파에서의 질량 의존 감쇠 진동 편차(부록 G)로부터의 독립적인 제약과 \textit{동시에} 적합시키는 데에 있다. 이 모든 데이터셋이 공동으로 고려될 때, 보편적 지수 $\beta=0.67$에 대한 사례는 설득력이 있게 된다.
		
		\subsection{향후 실험에 대한 예측}
		
		$\sim 1\,\mu$K-arcmin 잡음의 예상 감도를 가진 CMB‑S4는, YUCT 매개변수가 최적 적합 값에 가까우면 회전 $B$-모드를 $>5\sigma$로 검출할 수 있을 것이다. LiteBIRD는 더 큰 각도 규모에서의 결정적인 교차 확인을 제공할 것이다. 더욱이, 특정한 비가우스성 패턴($f_{\mathrm{NL}}^{\mathrm{local}} \approx 1.1$)은 향후 대규모 구조 서베이의 도달 범위 내에 있다.
		
		\section{논의와 결론}
		
		우리는 야쿠셰프 통합 조정 이론 내에서 소규모 CMB 비등방성의 최초의 완전한 취급을 제시하였다. 핵심 결과는:
		\begin{enumerate}
			\item \textbf{수정된 볼츠만 방정식}: 온도 의존 중력과 조정 점성을 통합하며, YUCT 섹터 간 결합($\kappa_{01}$)으로부터 유도됨.
			\item \textbf{YUCT 대수적 루프로의 연결}: 새로운 매개변수 $\varepsilon_0$과 $\eta$가 임의적이지 않고 보편적 상수 $\beta=2/3$, $\kappa_c=1/3$, $S_{\mathrm{odd}}$, $S_{\mathrm{even}}$에 연결되어 있음을 보여줌. 이로써 유효 자유 매개변수의 수가 줄어들고 물리적 동기가 강화됨.
			\item \textbf{Planck 높은 $\ell$ 데이터에 대한 개선된 적합}: 두 개의 추가 매개변수에 대해 $\Delta\chi^2 = -7.6$. $\ell \sim 2000$--$2500$에서의 초과 파워는 조정 점성으로 인한 실크 감쇠의 감소에 의해 자연스럽게 설명됨.
			\item \textbf{회전 $B$-모드의 독특한 예측}: 특징적인 $\ell^{-2}$ 형태를 가지며, 향후 CMB 실험으로 검증 가능하고, 원시 비가우스성의 특정 패턴도 예측함.
		\end{enumerate}
		
		우리는 현재 분석의 한계를 논의하였다: 통계적 개선은 단독으로 볼 때 미미하며, 표준 매개변수와의 축퇴가 존재한다. Planck 데이터의 계통적 효과도 높은 $\ell$ 잔차에 기여할 수 있다. 그럼에도 불구하고, YUCT 확장은 보편적 지수 $\beta = 0.67$과 대수적 루프를 통해 초기 우주를 실험실 물리학에 연결하는 일관된 물리적 묘사를 제공한다. 이 틀의 진정한 강점은 단일한 보편적 상수 집합으로부터, 이질적인 데이터셋(CMB, 백색왜성, 중력파)에 걸친 다수의 상관된 예측을 할 수 있는 능력에 있다.
		
		향후 연구는 완전한 19차원 라그랑지안으로부터의 조정 점성의 더 상세한 계산, 그리고 CMB, 대규모 구조, 천체물리학적 데이터를 결합한 공동 MCMC 분석을 포함해야 한다. 여기서 개발된 방법론은 대규모 구조 관측량에도 확장될 수 있어, 모든 우주론적 규모에 걸친 조정 패러다임의 포괄적인 테스트를 제공할 것이다.
		
		\section*{감사의 글}
		
		저자는 CAMB와 MontePython의 개발자들이 그들의 코드를 공개한 것에 감사한다.
		
		\section*{데이터 이용 가능성}
		
		수정된 CAMB 코드와 MCMC 연쇄는 YPSDC 리포지토리 \url{https://github.com/Alexey-Yakushev-YUCT/YPSDC}에서 이용 가능하다.
		
		\begin{thebibliography}{99}
			\bibitem{Planck2018} Planck Collaboration, \emph{Astron. Astrophys.} \textbf{641}, A6 (2020).
			\bibitem{Ma1995} C.-P. Ma, E. Bertschinger, \emph{Astrophys. J.} \textbf{455}, 7 (1995).
			\bibitem{Silk1968} J. Silk, \emph{Astrophys. J.} \textbf{151}, 459 (1968).
			\bibitem{Dodelson2003} S. Dodelson, \emph{Modern Cosmology}, Academic Press (2003).
			\bibitem{Hu1995} W. Hu, N. Sugiyama, \emph{Astrophys. J.} \textbf{444}, 489 (1995).
			\bibitem{Cooray2005} A. Cooray, D. Baumann, \emph{Phys. Rev. D} \textbf{71}, 063002 (2005).
			\bibitem{Lewis2000} A. Lewis, A. Challinor, A. Lasenby, \emph{Astrophys. J.} \textbf{538}, 473 (2000).
			\bibitem{Audren2013} B. Audren et al., \emph{J. Cosmol. Astropart. Phys.} \textbf{1302}, 001 (2013).
			\bibitem{Uzan2011} J.-P. Uzan, \emph{Living Rev. Relativ.} \textbf{14}, 2 (2011).
			\bibitem{BICEP2021} BICEP/Keck Collaboration, \emph{Phys. Rev. Lett.} \textbf{127}, 151301 (2021).
			\bibitem{AppendixP} A.V. Yakushev, ``Appendix P: Temperature Dependence of Gravity,'' in \emph{YUCT}, Zenodo (2026).
			\bibitem{AppendixM} A.V. Yakushev, ``Appendix M: YUCT Cosmology — Inflation as a Coordination Phase Transition,'' in \emph{YUCT}, Zenodo (2026).
			\bibitem{AppendixΩ} A.V. Yakushev, ``Appendix $\Omega$: The Global Rotation of the Universe and Its New Minimum Age,'' in \emph{YUCT}, Zenodo (2026).
			\bibitem{AppendixA} A.V. Yakushev, ``Appendix A: Practical Guide to Using YUCT V35.0,'' in \emph{YUCT}, Zenodo (2026).
			\bibitem{AppendixY} A.V. Yakushev, ``Appendix Y: Mathematical Foundations of YUCT,'' in \emph{YUCT}, Zenodo (2026).
			\bibitem{AppendixLambda} A.V. Yakushev, ``Appendix $\Lambda$: Resolution of the Hierarchy and Cosmological Constant Problems within YUCT,'' in \emph{YUCT}, Zenodo (2026).
			\bibitem{AppendixFerra} A.V. Yakushev, ``Appendix Ferra1.2: Universal Coordination Constants for Ordered and Disordered Phases,'' in \emph{YUCT}, Zenodo (2026).
		\end{thebibliography}
		
	\end{multicols}
\end{document}